\chapter{La valutazione e l'allenamento del metabolismo aerobico}
\label{cap:metabolismoaerobico}

% ---------------------------------------------------------------------------
\section{I test più utilizzati}

I test di valutazione del metabolismo aerobico più utilizzati nel futsal sono due famiglie:

\begin{enumerate}
  \item il \textbf{test Yo-Yo}, nelle versioni \textit{Endurance} e \textit{Intermittent Recovery};
  \item il \textbf{FIET} (\textit{Futsal Intermittent Endurance Test}).
\end{enumerate}

\rubrica{Perché proprio questi due}
Entrambi sono \textbf{test specifici} per il calcio a 5:

\begin{itemize}
  \item lo \textbf{Yo-Yo} --- e molto più l'\textit{intermittent recovery} che l'\textit{endurance}
    --- si avvicina al \textbf{modello di prestazione}, perché prevede un \textbf{cambio di
    direzione} e sollecita quindi le componenti di \textbf{accelerazione} e \textbf{decelerazione};
    è inoltre un test \textbf{incrementale}, in cui si raggiungono frequenze cardiache molto
    elevate;
  \item il \textbf{FIET} è ancora più specifico per il futsal: prevede una percorrenza a navetta di
    \textbf{3 $\times$ 15\,m} seguita da un tempo di recupero;
  \item hanno \textbf{durate confrontabili}, sono \textbf{facili da strutturare in campo} e
    richiedono \textbf{pochissimo materiale};
  \item danno un'indicazione sulla \textbf{fitness aerobica} del giocatore.
\end{itemize}

% ---------------------------------------------------------------------------
\section{Il test Yo-Yo: impianto generale}

\textbf{Yo-Yo}: test che si svolgono eseguendo attività di \textbf{corsa a navetta su base di
20\,m}. Sono giunti in Italia nel \textbf{1997} (\textbf{Bangsbo}, 1997) e vengono impiegati
principalmente negli \textbf{sport di squadra}.

\rubrica{Le tre versioni}
\begin{enumerate}
  \item \textbf{Yo-Yo Endurance Test};
  \item \textbf{Yo-Yo Intermittent Endurance Test};
  \item \textbf{Yo-Yo Intermittent Recovery Test}.
\end{enumerate}

Le varie versioni esistono per rendere il test fruibile a operatori di discipline diverse.

\rubrica{La lepre acustica}
Per ognuna delle tre versioni esistono \textbf{tracce audio} con \textit{beep} sonori che dettano
la velocità di percorrenza:

\begin{itemize}
  \item la velocità cresce gradualmente e deve essere mantenuta con \textbf{continuità}: a ogni
    \textit{beep} il soggetto si deve trovare al rispettivo delimitatore;
  \item essendo la \textbf{distanza fissa}, all'aumentare della \textbf{frequenza del beep} aumenta
    necessariamente la \textbf{velocità di esecuzione};
  \item il \textbf{cambio di direzione} va eseguito proprio in coincidenza del \textit{beep}
    acustico.
\end{itemize}

\rubrica{I due livelli}
Ognuna delle tre versioni ha \textbf{due livelli di esecuzione}:

\begin{itemize}
  \item \textbf{livello 1}: per chi è poco o per nulla allenato;
  \item \textbf{livello 2}: per chi è in un buono stato di allenamento, o per chi è riuscito a
    completare il livello 1;
  \item la differenza consiste nel \textbf{partire, nel livello 2, da una velocità di corsa
    maggiore} (Bangsbo, 1997).
\end{itemize}

\rubrica{Dallo stage alla velocità reale}
I numeri degli \textbf{step} indicati dal protocollo \textbf{non corrispondono alle velocità
reali}, ma indicano soltanto lo \textbf{stadio raggiunto}. Il prof.\ \textbf{Toschi} (2002) ha
individuato la formula di conversione in velocità reale, in km/h:

\[
  v = 7,5 + \frac{n^\circ\ stage}{2}
\]

\subsection{Consigli per l'esecuzione}

\rubrica{Prima di partire}
\begin{itemize}
  \item chiedere ai giocatori se hanno \textbf{già eseguito il test} in carriera e, in caso
    contrario, spiegare il protocollo nella maniera più chiara possibile;
  \item \textbf{cambiare piede di appoggio a ogni cambio di senso}, per evitare di dare
    sovraccarico funzionale a una sola parte del corpo;
  \item far \textbf{concentrare il giocatore sul beep sonoro}, così da mantenere la corretta
    andatura di corsa.
\end{itemize}

\rubrica{L'errore tipico: anticipare il beep}
I livelli 1 --- sia dell'\textit{endurance} sia del \textit{recovery} --- partono a velocità di
percorrenza \textbf{basse}, e la tendenza dei giocatori è di \textbf{sovrastimarle}: arrivano ai
20\,m \textbf{anticipando} il \textit{beep}. In quel caso l'indicazione è \textbf{fermarsi dietro
la linea}, attendere il \textit{beep} e solo allora ripartire.

\rubrica{Gli operatori}
Quando si valutano più soggetti contemporaneamente è consigliabile la presenza di \textbf{almeno
due operatori}, ognuno posto all'altezza di un delimitatore. Il motivo sta nei \textbf{due modi in
cui un giocatore termina il test}:

\begin{enumerate}
  \item su \textbf{base volontaria}: ha esaurito le energie e si ferma;
  \item perché, \textbf{pur volendo continuare}, non riesce più ad arrivare ai 20\,m in coincidenza
    con il \textit{beep} --- generalmente si attendono \textbf{due o al massimo tre step} prima di
    fermarlo.
\end{enumerate}

\rubrica{Testare atleti che non si conoscono}
L'indicazione di stop deve essere \textbf{precisa e individuale}: un generico ``fermati'' verrebbe
frainteso dai giocatori vicini. Due soluzioni:

\begin{itemize}
  \item \textbf{farsi affiancare da almeno una persona che li conosca};
  \item \textbf{numerare le corsie}, così che a ogni giocatore corrisponda un numero e lo stop si
    chiami per numero di corsia.
\end{itemize}

% ---------------------------------------------------------------------------
\section{Lo Yo-Yo Endurance Test}

Versione aggiornata del \textbf{test di Léger} e del \textbf{Multistage Fitness Test}. È un test
\textbf{continuo} e valuta la capacità di eseguire esercizio fisico in maniera \textbf{continuativa
e prolungata}.

\rubrica{Il protocollo}
\begin{itemize}
  \item due delimitatori posti a \textbf{20\,m} l'uno dall'altro, con una \textbf{corsia di 2\,m};
  \item il soggetto corre \textbf{avanti e indietro} a velocità progressivamente crescenti, pari a
    \textbf{$+$0,5\,km/h ogni minuto circa}, dettate dalla lepre acustica;
  \item la cadenza sonora prevede un \textit{beep} \textbf{alla partenza}, uno \textbf{al cambio di
    direzione} e uno \textbf{nuovamente in posizione di partenza};
  \item il test \textbf{termina} quando il soggetto non è più in grado di mantenere il ritmo
    imposto;
  \item \textbf{durata}: da \textbf{5} a \textbf{20 minuti}.
\end{itemize}

\rubrica{Che cosa si registra}
Al termine si registrano:

\begin{itemize}
  \item lo \textbf{step di fermata};
  \item le \textbf{navette} in esso coperte;
  \item il \textbf{numero di metri percorsi} --- direttamente calcolabile dal numero di navette,
    essendo nota la distanza;
  \item la \textbf{velocità finale}.
\end{itemize}

\rubrica{Il $VO_2max$ teorico}
Dai risultati si possono ricavare valori di \textbf{$VO_2max$ teorico} tramite una \textbf{tabella
di conversione} valida per ogni fascia di età (Castagna, 1999), scaturita da ricerche scientifiche
(Bangsbo, 1997).

\begin{itemize}
  \item la stima \textbf{presuppone il massimo impegno} del soggetto: è questo il \textbf{limite
    più grande} del test e, in generale, di \textbf{tutti i test che prevedono un esaurimento}.
\end{itemize}

\rubrica{Il dato utile per l'allenamento}
Il riferimento da portarsi a casa è la \textbf{velocità finale}, che serve a impostare un
\textbf{allenamento continuo}. Nel calcio a 5 lo si somministra \textbf{raramente}, perché è
estremamente lontano dal modello di prestazione: trova spazio in fasi di \textbf{mantenimento della
forma} o di \textbf{recupero da un infortunio}.

% ---------------------------------------------------------------------------
\section{Lo Yo-Yo Intermittent Recovery Test}

Valuta la \textbf{capacità di recupero} del soggetto durante uno sforzo progressivamente crescente:
non solo la capacità di \textbf{produrre energia} per via aerobica, ma anche quella di
\textbf{recuperare velocemente} tra uno sforzo e l'altro.

\rubrica{Il protocollo}
\begin{itemize}
  \item corsa a navetta tra due delimitatori distanti \textbf{20\,m}, su una \textbf{corsia di
    2\,m};
  \item al termine di ogni \textbf{frazione di 40\,m} (20 $+$ 20\,m) si eseguono \textbf{10\,s di
    recupero attivo} (\textit{jogging}), girando dietro un \textbf{terzo delimitatore posto a 5\,m}
    dietro e leggermente a fianco della linea di partenza --- una corsa molto lenta, di scarico,
    fino a tornare fermi sulla linea di partenza;
  \item la sequenza sonora è: \textbf{1\textsuperscript{o} \textit{beep}} alla partenza,
    \textbf{2\textsuperscript{o}} al delimitatore dei 20\,m, \textbf{3\textsuperscript{o}} al
    rientro sulla linea di partenza;
  \item sia la \textbf{velocità di corsa} sia il \textbf{tempo fisso di recupero di 10\,s} sono
    dettati per tutta la prova da un \textit{beep} pre-registrato;
  \item se il giocatore \textbf{anticipa} il rientro, si ferma e attende il \textit{beep};
  \item \textbf{durata}: da \textbf{2} a \textbf{15 minuti} --- i 2 minuti sono un caso limite, e in
    campo il test non supera mediamente i 15--20 minuti.
\end{itemize}

\rubrica{Le velocità iniziali}
\begin{itemize}
  \item \textbf{livello 1}: \textbf{10\,km/h} (step 5);
  \item \textbf{livello 2}: \textbf{13\,km/h} (step 11) (Bangsbo, 1997).
\end{itemize}

Nel livello 1 l'indicazione è di \textbf{non partire sprintando}: a 10\,km/h la prestazione è una
corsa a \textbf{moderata intensità}.

\rubrica{Il criterio di arresto}
\begin{itemize}
  \item \textbf{prima volta} che il soggetto non raggiunge la linea di partenza in concomitanza con
    il \textit{beep}: \textbf{ammonizione};
  \item \textbf{seconda volta}: il soggetto viene \textbf{fermato}.
\end{itemize}

\rubrica{Che cosa si registra}
\begin{itemize}
  \item gli \textbf{step};
  \item il \textbf{numero di navette}, compresa l'ultima;
  \item il \textbf{numero di metri percorsi};
  \item la \textbf{velocità finale}.
\end{itemize}

\rubrica{Il dato utile per l'allenamento}
Anche qui la \textbf{velocità finale} dà un'indicazione diretta, questa volta per l'allenamento a
\textbf{regime intermittente} con modalità di esecuzione \textbf{a navetta}.

% ---------------------------------------------------------------------------
\section{Capacità aerobica e Yo-Yo IR1 nel futsal}

Studio di \textbf{Boullosa e coll.} (2013), \textit{Relationship between Aerobic Capacity and
Yo-Yo IR1 Performance in Brazilian Professional Futsal Players}, \textit{Asian J Sports Med}
4(3): 230--234.

\rubrica{Scopo}
Valutare le capacità aerobiche, mettendo in \textbf{relazione il $VO_2$ e il risultato dello Yo-Yo
IR test} in giocatori di futsal.

\rubrica{Campione}
$n = 15$ giocatori professionisti brasiliani, di livello di qualificazione elevato.

\begin{table}[htbp]
  \centering
  \small
  \renewcommand{\arraystretch}{1.3}
  \begin{tabularx}{\textwidth}{|X|c|}
    \hline
    \textbf{Parametro} & \textbf{Media (DS)} \\
    \hline
    Età (anni) & 25,9 (5,1) \\
    \hline
    Massa corporea (kg) & 74,37 (6,02) \\
    \hline
    Statura (m) & 1,77 (0,04) \\
    \hline
    Yo-Yo Intermittent Recovery Test livello 1 (m) & 1226 (282) \\
    \hline
    $FC$ massima registrata nel test Yo-Yo (bpm) & 184 (15) \\
    \hline
    Lattato massimo dopo il test Yo-Yo (mmol/L) & 8,17 (1,63) \\
    \hline
    Velocità aerobica massima nel \textit{ramp test} (km/h) & 17,69 (1,88) \\
    \hline
    $FC$ massima registrata nel \textit{ramp test} (bpm) & 184 (5) \\
    \hline
    $VO_2max$ (ml$\cdot$kg$^{-1}\cdot$min$^{-1}$) & 57,25 (6,35) \\
    \hline
    $VO_2max$ (ml$\cdot$kg$^{-0,75}\cdot$min$^{-1}$) & 76,34 (8,46) \\
    \hline
  \end{tabularx}
\end{table}

\rubrica{Una nota sulla frequenza cardiaca}
I 184 bpm registrati nello Yo-Yo \textbf{non sono la $FC_{max}$} dei giocatori: nel test
intermittente intervengono a fermare la prestazione le componenti del \textbf{metabolismo
anaerobico lattacido}, che fanno insorgere una fatica momentanea e portano il giocatore a
terminare la prova prima. Per questo lo \textbf{Yo-Yo IR test difficilmente viene utilizzato per
stabilire la frequenza cardiaca massima}.

\rubrica{I risultati}
\begin{itemize}
  \item \textbf{$VO_2max$ e risultato del test}: \textbf{nessuna relazione} ($r = 0,098$;
    $p = 0,719$) con il $VO_2max$ rilevato durante il test incrementale su \textit{treadmill};
  \item \textbf{velocità di picco e risultato del test}: correlazione \textbf{statisticamente
    significativa} ($r = 0,641$; $p = 0,007$) tra la \textbf{$S_{peak}$} raggiunta nel test su
    \textit{treadmill} e il risultato dello Yo-Yo.
\end{itemize}

\rubrica{\textit{Take home message}}
\begin{itemize}
  \item la \textbf{velocità di picco} presenta una correlazione con il risultato nello Yo-Yo IR
    livello 1, \textbf{ma non il $VO_2max$};
  \item il fatto che il $VO_2max$ non sia valutabile con lo Yo-Yo IR è \textbf{connaturato
    all'attività intermittente} del test: resta un test interessante per l'analisi dell'\textbf{alta
    intensità};
  \item per strutturare un \textbf{profilo della fitness aerobica} del giocatore di futsal è
    necessario svolgere \textbf{entrambi}:
    \begin{itemize}
      \item un \textbf{test continuo}, per il $VO_2max$ e la \textbf{potenza aerobica};
      \item un \textbf{test intermittente ad alta intensità};
    \end{itemize}
  \item il test continuo può essere lo \textbf{Yo-Yo Endurance Test}, ma per avere la certezza del
    valore di massima potenza aerobica erogata è preferibile la \textbf{valutazione del $VO_2$ su
    \textit{treadmill}}.
\end{itemize}
